\BourbakiExerciseGroup

\begin{enumerate}
\def\labelenumi{\arabic{enumi})}
\item
  设 E 是交换域 K 上的一个结合、交换且保单位的代数；设 \(\mathbf{x} = (x_i)_{i \in I}\) 是 E 中两两可交换的元素族。设 U 是 \(\mathrm{K}(\mathrm{X}_i)_{i \in I}\) 的子环，由所有 \(f \in \mathrm{K}(\mathrm{X}_i)_{i \in I}\) （x 可代入 f）组成。证明：若 E 中非可逆元组成的集合是一个理想，则 U 中亦然。并证明此时 U 中非可逆元组成的理想是极大的。
\item
  \begin{enumerate}
  \def\labelenumii{\alph{enumii})}
  \tightlist
  \item
    设 K 为一个交换域，且 \(u = a_{m}X^{m} + a_{m+1}X^{m+1} + \cdots + a_{n}X^{n}\) 为 K{[}X{]} 中的多项式，使得 \(a_{m} \neq 0\) , \(a_{n} \neq 0\) ( \(0 \leq m \leq n\) )。证明：若 g 是 K(X) 中次数为 \(d \neq 0\) 的有理分式，则 \(u(g)\) 非零，且当 d \textgreater{} 0 时次数为 nd，当 d \textless{} 0 时次数为 md。
  \end{enumerate}
\end{enumerate}

\begin{enumerate}
\def\labelenumi{\alph{enumi})}
\setcounter{enumi}{1}
\tightlist
\item
  由此推断：K(X) 的每个非常数有理分式 g 都可代入 K(X) 中的每个有理分式（注意，若 g 的次数为 0，则存在 \(\alpha \in K\) 使得 \(g - \alpha\) 的次数 \textless{} 0）。
\end{enumerate}

\begin{enumerate}
\def\labelenumi{\arabic{enumi})}
\setcounter{enumi}{2}
\tightlist
\item
  一个有理分式 \(f \in K(X_1, X_2, \ldots, X_n)\) 被称为齐次的，如果它等于两个齐次多项式之商（分母 \(\neq 0\) ). 证明：为使 \(f\) 是齐次的，必要且充分条件是
\end{enumerate}

\[
f \left(\mathrm{ZX} _ {1}, \mathrm{ZX} _ {2}, \dots , \mathrm{ZX} _ {n}\right) = \mathrm{Z} ^ {d} f \left(\mathrm{X} _ {1}, \dots , \mathrm{X} _ {n}\right)
\]

其中 \(d\) 是 \(f\) 的次数。

\begin{enumerate}
\def\labelenumi{\arabic{enumi})}
\setcounter{enumi}{3}
\tightlist
\item
  证明：拉格朗日插值公式（IV，第16页，命题6）可以写成
\end{enumerate}

\[
f (\mathbf {X}) = \omega (\mathbf {X}) \sum_ {i = 1} ^ {n} \frac {\beta_ {i}}{\omega^ {\prime} (\alpha_ {i}) (\mathbf {X} - \alpha_ {i})}
\]

其中 \(\omega(X) = (X - \alpha_1)(X - \alpha_2) \ldots (X - \alpha_n)\) 。由此推出，如果 \(g\) 是 \(K[X]\) 中次数 \(\leqslant n - 2\) 的多项式，则

\[
\sum_ {i = 1} ^ {n} \frac {g (\alpha_ {i})}{\omega^ {\prime} (\alpha_ {i})} = 0
\]

（考虑多项式 \(f(\mathbf{X}) = \mathbf{X}g(\mathbf{X})\) ）。

5) 设 \(\mathbf{K}\) 为特征 0 的交换域。证明：若有理分式 \(u \in \mathbf{K}(\mathbf{X})\) 使得 \(\mathbf{D}u = 0\) ，则 \(u\) 是常数。

\begin{enumerate}
\def\labelenumi{\arabic{enumi})}
\setcounter{enumi}{5}
\item
  将欧拉恒等式推广到特征 0 的域上的齐次有理分式（习题3）（考虑有理分式 \(\frac{1}{Z^m} f(ZX_1, ZX_2, ..., ZX_n)\) 关于 \(Z\) ，并利用习题5）。
\item
  设 \(\mathbf{K}\) 为一个交换域。证明：
\end{enumerate}

\[
[ \mathbf {K} (\mathrm{T}): \mathbf {K} ] = \operatorname{Sup} (\operatorname{Card} (\mathbf {K}), \operatorname{Card} (\mathbf {N})).
\]

(利用 \(1 / (T - a)\) （ \(a \in \mathbf{K}\) ）的线性无关性，来表明

\[
[ \mathrm{K} (\mathrm{T}): \mathrm{K} ] \geqslant \operatorname{Card} (\mathrm{K}).)
\]
% semantic-fragments: legacy-input-seam
\relax{}
